\documentclass[conference,a4paper]{IEEEtran}
\IEEEoverridecommandlockouts
\usepackage{iftex}
\ifLuaTeX
  \usepackage{fontspec}
\else
  \usepackage[utf8]{inputenc}
  \usepackage[T1]{fontenc}
\fi
\usepackage[spanish,es-noshorthands,es-tabla]{babel}
\usepackage{amsmath,amssymb}
\usepackage{graphicx}
\usepackage{booktabs}
\usepackage{array}
\usepackage{multirow}
\usepackage{xcolor}
\usepackage{tikz}
\usetikzlibrary{arrows.meta,positioning,fit}
\usepackage[hidelinks]{hyperref}
\usepackage{cite}
\decimalpoint
\newcommand{\pmsd}[2]{#1\,$\pm$\,#2}

\begin{document}

\title{Detección de enfermedades de la papa\\
con aprendizaje profundo:\\
comparación de ResNet-50 y Vision Transformer\\
con imágenes de campo}

\author{\IEEEauthorblockN{Gabriel Ángel Nizama Duque, Martín Alexis Valdiviezo Palma,\\ Juan José Pino Reyes, César Luis Jiménez Miranda}
\IEEEauthorblockA{Maestría en Inteligencia Artificial, Unidad de Posgrado FIIS\\
Universidad Nacional de Ingeniería, Lima, Perú\\
MIA-203 Redes Neuronales y Aprendizaje Profundo --- Equipo 4, Sección A}}

\maketitle

\begin{abstract}
Diagnosticar a tiempo las enfermedades de la papa es clave para un cultivo central en la seguridad alimentaria del Perú. Los modelos de visión suelen entrenarse con fotos de laboratorio, pero en campo hay fondos variados, manos, sombras y plásticos. Estudiamos si, con solo 2\,112 imágenes de entrenamiento tomadas en campo y 7 clases desbalanceadas, conviene una red convolucional (ResNet-50) o un Vision Transformer (ViT-S/16), y cuánto aporta el preentrenamiento en ImageNet. Primero auditamos el conjunto de datos público: encontramos 55 pares de duplicados exactos (6 de ellos con etiquetas contradictorias) y casi-duplicados, y construimos una partición 70/15/15 por grupos sin fuga de datos. Comparamos 6 configuraciones, con 3 semillas cada una: ResNet-50 y ViT-S/16, preentrenados y desde cero, y dos líneas base, un Random Forest con descriptores de color y textura y una CNN pequeña. Preentrenados, ResNet-50 y ViT-S/16 alcanzan macro-F1 de \pmsd{0.843}{0.012} y \pmsd{0.853}{0.013} en prueba, sin diferencia significativa según la prueba de McNemar ($p\geq0.16$). Ambos superan ampliamente al Random Forest (\pmsd{0.464}{0.021}). Sin preentrenamiento, el ViT cae a \pmsd{0.499}{0.028}, por debajo de la ResNet-50 entrenada desde cero (\pmsd{0.642}{0.007}). Es decir, con pocos datos el preentrenamiento importa más que la arquitectura. En los ejemplos revisados con Grad-CAM, la ResNet enfoca las regiones con lesión, mientras que la CNN pequeña a veces usa atajos del fondo. Publicamos el código, la partición, los pesos y una demo reproducible.
\end{abstract}

\begin{IEEEkeywords}
aprendizaje profundo, Vision Transformer, ResNet, transferencia de aprendizaje, enfermedades de la papa, Grad-CAM, pocos datos.
\end{IEEEkeywords}

%=====================================================================
\section{Introducción}
La papa es uno de los cultivos más importantes del Perú y del mundo. Las enfermedades foliares, como la rancha (\textit{Phytophthora infestans}), las virosis o el daño por insectos, reducen el rendimiento si no se detectan a tiempo. El diagnóstico visual depende de especialistas que no siempre están disponibles en campo. Por eso, un clasificador de imágenes ejecutable desde un celular podría servir como apoyo temprano.

El aprendizaje profundo ya logra exactitudes casi perfectas en conjuntos de laboratorio como PlantVillage~\cite{mohanty2016plant,ferentinos2018plant}, con hojas aisladas sobre fondo uniforme. Sin embargo, esos modelos se degradan con fotos de campo~\cite{mohanty2016plant}. El conjunto de Shabrina \textit{et al.}~\cite{shabrina2024potato} reúne 3\,076 fotos de hojas de papa tomadas en campo en Java Central (Indonesia), en 7 clases. Sus líneas base reportan exactitudes de solo 0.68--0.74.

Al mismo tiempo, los Vision Transformers (ViT)~\cite{dosovitskiy2021vit} han desplazado a las CNN en muchos problemas. Sin embargo, se sabe que necesitan muchos datos o un preentrenamiento fuerte~\cite{dosovitskiy2021vit,touvron2021deit,steiner2022howtotrain}. Esto plantea una pregunta práctica para equipos con pocos datos y poco cómputo: \emph{¿qué arquitectura elegir y cuánto depende el resultado del preentrenamiento?}

\textbf{Contribuciones.}
\begin{enumerate}
\item \textbf{Auditoría y partición sin fuga.} Auditamos el conjunto de datos público: 110 imágenes en 55 pares de duplicados exactos, 12 de ellas con etiquetas contradictorias, y 64 imágenes en grupos de casi-duplicados. Construimos una partición 70/15/15 por grupos, reproducible con un comando.
\item \textbf{Comparación controlada.} Comparamos ResNet-50 y ViT-S/16 (preentrenados y desde cero), una CNN pequeña y un Random Forest, con 3 semillas, métricas adecuadas al desbalance (macro-F1) y prueba de significancia (McNemar).
\item \textbf{Ablación del preentrenamiento.} Cuantificamos que aporta $+0.355$ de macro-F1 al ViT y $+0.201$ a la ResNet.
\item \textbf{Interpretabilidad y demo.} Analizamos los errores con Grad-CAM, incluyendo una medida de concentración de los mapas, y publicamos el modelo final en una app web que lo ejecuta en el celular, además de una demo reproducible en Colab.
\end{enumerate}

El resto del artículo se organiza así: la Sección~\ref{sec:related} revisa trabajos relacionados; la Sección~\ref{sec:method} describe la metodología; la Sección~\ref{sec:results} presenta los resultados, que se discuten en la Sección~\ref{sec:discussion}; la Sección~\ref{sec:conclusions} concluye.

%=====================================================================
\section{Trabajos relacionados}\label{sec:related}
\textbf{CNN y transferencia de aprendizaje.} Las redes residuales~\cite{he2016resnet} permitieron entrenar CNN profundas mediante conexiones de identidad. Su preentrenamiento en ImageNet~\cite{deng2009imagenet} es el punto de partida estándar con pocos datos. En plantas, Mohanty \textit{et al.}~\cite{mohanty2016plant} y Ferentinos~\cite{ferentinos2018plant} alcanzaron más de 99\,\% de exactitud en PlantVillage, pero advierten que el desempeño cae con imágenes de campo. Nuestro trabajo usa precisamente fotos de campo.

\textbf{Transformers en visión.} El Transformer~\cite{vaswani2017attention} reemplaza la convolución por autoatención. El ViT~\cite{dosovitskiy2021vit} lo aplica a parches de $16\times16$ píxeles y supera a las CNN solo cuando se preentrena en conjuntos muy grandes. DeiT~\cite{touvron2021deit} y Steiner \textit{et al.}~\cite{steiner2022howtotrain} mostraron que la regularización y el preentrenamiento compensan en parte la falta de sesgo inductivo. Raghu \textit{et al.}~\cite{raghu2021vitsee} encontraron que el ViT integra información global desde las primeras capas, a diferencia de la CNN. Nosotros medimos este contraste en un problema real con solo $\sim$2\,100 imágenes.

\textbf{Papa en condiciones no controladas.} Shabrina \textit{et al.}~\cite{shabrina2024potato} publicaron el conjunto que usamos y reportaron líneas base con una sola partición 90/10: ResNet-50 obtiene 0.682 de exactitud y EfficientNetV2-B3, 0.736. Sangar y Rajasekar~\cite{sangar2025potato} combinan EfficientNet-LITE con un SVM en entornos diversos. Ninguno de los dos trabajos controla duplicados entre particiones, reporta varias semillas ni compara con Transformers.

\textbf{Interpretabilidad y atajos.} Grad-CAM~\cite{selvaraju2017gradcam} localiza las regiones que más influyen en una clase. Su aplicación a ViT es menos fiel que en CNN~\cite{chefer2021transformer}, lo que consideramos al interpretar. Geirhos \textit{et al.}~\cite{geirhos2020shortcut} describen el \emph{aprendizaje por atajos}: decidir por pistas espurias del fondo. Lo observamos en nuestra CNN pequeña.

%=====================================================================
\section{Metodología}\label{sec:method}
\begin{figure*}[t]
\centering
\resizebox{\textwidth}{!}{%
\begin{tikzpicture}[node distance=4mm, font=\footnotesize,
  caja/.style={draw, rounded corners=2pt, align=center, minimum height=10mm, inner sep=3pt, fill=#1},
  caja/.default=white, flecha/.style={-{Latex[length=2mm]}, thick}]
\node[caja=gray!10] (raw) {3\,076 fotos\\1500$\times$1500\\7 clases};
\node[caja=gray!10, right=of raw] (dedup) {Deduplicación\\sha256 (exactos)\\dHash $\leq4$ (casi)};
\node[caja=gray!10, right=of dedup] (split) {Partición por grupos\\70/15/15, semilla 42\\2112 / 450 / 453};
\node[caja=gray!10, right=of split] (aug) {256$\to$224 px\\aumento: recorte,\\volteo, rotación, color};
\node[caja=blue!8, right=10mm of aug, yshift=9mm] (cnn) {ResNet-50 / CNN pequeña\\(convolución local)};
\node[caja=orange!12, right=10mm of aug, yshift=-9mm] (vit) {ViT-S/16: 196 parches\\+ autoatención global};
\node[caja=green!8, right=10mm of cnn, yshift=-9mm] (eval) {Mejor época por\\macro-F1 de validación\\$\Rightarrow$ prueba (1 vez)};
\node[caja=gray!10, right=of eval] (xai) {Grad-CAM,\\McNemar,\\demo};
\draw[flecha] (raw) -- (dedup); \draw[flecha] (dedup) -- (split); \draw[flecha] (split) -- (aug);
\draw[flecha] (aug.east) -- (cnn.west); \draw[flecha] (aug.east) -- (vit.west);
\draw[flecha] (cnn.east) -- (eval.west); \draw[flecha] (vit.east) -- (eval.west); \draw[flecha] (eval) -- (xai);
\end{tikzpicture}}
\caption{Flujo del estudio. La partición se hace por grupos de casi-duplicados para evitar fuga de datos; el conjunto de prueba solo se usa una vez, al final de cada corrida.}
\label{fig:pipeline}
\end{figure*}

La Fig.~\ref{fig:pipeline} resume el flujo del estudio: auditoría y limpieza de los datos, partición sin fuga, entrenamiento de los modelos y evaluación única en prueba.

\subsection{Datos}
Usamos el \textit{Potato Leaf Disease Dataset in Uncontrolled Environment}~\cite{shabrina2024potato,shabrina2023mendeley} (Mendeley Data, licencia CC BY 4.0). Contiene 3\,076 fotos JPG de $1500\times1500$ píxeles tomadas en campo. Las 7 clases son bacteria, hongos (\textit{fungi}), sana, nematodos, plaga de insectos (\textit{pest}), \textit{Phytophthora} y virus. El desbalance es marcado: de 67 imágenes (nematodos) a 740 (hongos) tras la limpieza. La Fig.~\ref{fig:dataset} muestra un ejemplo por clase y la distribución por partición. Las fotos incluyen manos, plásticos de cobertura y varias hojas por imagen.

\begin{figure*}[t]
\centering
\includegraphics[width=0.92\textwidth]{figs/dataset.png}
\caption{Conjunto de datos. Arriba: una imagen de prueba por clase, tomada en campo con fondos no controlados (plástico, manos, otras plantas). Abajo: número de imágenes por clase y partición tras eliminar duplicados; la clase mayor (hongos) tiene 11 veces más imágenes que la menor (nematodos).}
\label{fig:dataset}
\end{figure*}

\textbf{Auditoría.} Calculamos el hash sha256 de cada archivo y encontramos 110 imágenes que forman 55 pares de duplicados exactos. En 6 pares, la misma foto aparece en dos clases distintas. Eliminamos esas 12 imágenes, por tener etiqueta contradictoria, y 49 copias redundantes, de modo que quedan 3\,015. Además, calculamos un \textit{hash} perceptual dHash de 64 bits sobre la versión de $256\times256$. Unimos con \textit{union-find} los pares con distancia de Hamming $\leq4$, lo que da 64 imágenes en grupos de casi-duplicados.

\textbf{Partición.} Hacemos una partición estratificada por grupo 70/15/15 (semilla 42): 2\,112 imágenes de entrenamiento, 450 de validación y 453 de prueba. Cada grupo de casi-duplicados cae completo en una sola partición, para que el modelo no se evalúe con variantes de fotos que ya vio. La validación se usa para elegir la mejor época y el modelo final; la prueba se evalúa una única vez por corrida.

\subsection{Modelos}
\begin{figure*}[t]
\centering
\resizebox{\textwidth}{!}{%
\begin{tikzpicture}[font=\footnotesize, node distance=3.5mm,
  b/.style={draw, rounded corners=2pt, align=center, minimum height=11mm, inner sep=2.5pt, fill=#1},
  b/.default=white, f/.style={-{Latex[length=1.6mm]}, thick},
  etq/.style={font=\footnotesize\bfseries, anchor=east, align=right, text width=17mm},
  dim/.style={font=\scriptsize, text=black!60, anchor=north}]
% --- CNN pequeña
\node[etq] (l1) at (0,0) {CNN\\pequeña\\(1.2\,M)};
\node[b=gray!10, right=2mm of l1] (c0) {Imagen\\$224{\times}224{\times}3$};
\node[b=blue!8, right=of c0] (c1) {Bloque 1\\{[Conv$3{\times}3$-BN-ReLU]}$\times2$\\+ MaxPool, 32};
\node[b=blue!8, right=of c1] (c2) {Bloque 2\\64};
\node[b=blue!8, right=of c2] (c3) {Bloque 3\\128};
\node[b=blue!8, right=of c3, line width=1.2pt] (c4) {Bloque 4\\256};
\node[b=gray!10, right=of c4] (c5) {Pooling\\global};
\node[b=green!10, right=of c5] (c6) {Dropout 0.3\\Lineal $\to$ 7};
\foreach \a/\bb in {c0/c1,c1/c2,c2/c3,c3/c4,c4/c5,c5/c6} \draw[f] (\a) -- (\bb);
\node[dim] at (c1.south) {$112^2$}; \node[dim] at (c2.south) {$56^2$};
\node[dim] at (c3.south) {$28^2$}; \node[dim] at (c4.south) {$14^2$};
% --- ResNet-50
\node[etq] (l2) at (0,-2.0) {ResNet-50\\(23.5\,M)};
\node[b=gray!10, right=2mm of l2] (r0) {Imagen\\$224{\times}224{\times}3$};
\node[b=blue!8, right=of r0] (r1) {Conv $7{\times}7$/2, 64\\+ MaxPool};
\node[b=blue!8, right=of r1] (r2) {Etapa 1\\3 bloques\\256};
\node[b=blue!8, right=of r2] (r3) {Etapa 2\\4 bloques\\512};
\node[b=blue!8, right=of r3] (r4) {Etapa 3\\6 bloques\\1024};
\node[b=blue!8, right=of r4, line width=1.2pt] (r5) {Etapa 4\\3 bloques\\2048};
\node[b=gray!10, right=of r5] (r6) {Pooling\\global};
\node[b=green!10, right=of r6] (r7) {Dropout 0.2\\Lineal $\to$ 7};
\foreach \a/\bb in {r0/r1,r1/r2,r2/r3,r3/r4,r4/r5,r5/r6,r6/r7} \draw[f] (\a) -- (\bb);
\node[dim] at (r1.south) {$56^2$}; \node[dim] at (r2.south) {$56^2$}; \node[dim] at (r3.south) {$28^2$};
\node[dim] at (r4.south) {$14^2$}; \node[dim] at (r5.south) {$7^2$};
\node[b=white, right=6mm of r7, font=\scriptsize, minimum height=9mm] (bott) {Bloque cuello de botella:\\$1{\times}1 \to 3{\times}3 \to 1{\times}1$\\$\mathbf{y}=\mathcal{F}(\mathbf{x})+\mathbf{x}$};
\draw[dashed, black!50] (r5.north east) -- (bott.north west);
% --- ViT-S/16
\node[etq] (l3) at (0,-4.0) {ViT-S/16\\(21.7\,M)};
\node[b=gray!10, right=2mm of l3] (v0) {Imagen\\$224{\times}224{\times}3$};
\node[b=orange!12, right=of v0] (v1) {196 parches\\$16{\times}16$};
\node[b=orange!12, right=of v1] (v2) {Proyección lineal\\$\to d{=}384$\\+ [CLS] + posición};
\node[b=orange!20, right=of v2, line width=1.2pt, text width=36mm] (v3) {Bloque Transformer $\times 12$:\\LN $\to$ MHSA (6 cabezas) $\to$ +\\LN $\to$ MLP (1536, GELU) $\to$ +};
\node[b=gray!10, right=of v3] (v4) {LayerNorm\\token [CLS]};
\node[b=green!10, right=of v4] (v5) {Dropout 0.2\\Lineal $\to$ 7};
\foreach \a/\bb in {v0/v1,v1/v2,v2/v3,v3/v4,v4/v5} \draw[f] (\a) -- (\bb);
\node[dim] at (v2.south) {$197{\times}384$}; \node[dim] at (v3.south) {$197{\times}384$};
\end{tikzpicture}}
\caption{Arquitecturas comparadas. Las cifras bajo cada caja son el tamaño espacial del mapa ($h{\times}w$) o de la secuencia de \textit{tokens}. El borde grueso marca la capa donde se calcula Grad-CAM. Los tres modelos terminan en una capa lineal de 7 salidas seguida de \textit{softmax}.}
\label{fig:arquitecturas}
\end{figure*}

La Fig.~\ref{fig:arquitecturas} resume las tres redes.

\textbf{Línea base clásica.} Un Random Forest~\cite{breiman2001rf} con 1\,834 descriptores por imagen:
\begin{itemize}
\item Histograma HSV de 48 dimensiones.
\item Medias y desviaciones RGB/HSV (12).
\item HOG~\cite{dalal2005hog} sobre $128\times128$ (1\,764).
\item LBP uniforme~\cite{ojala2002lbp} (10).
\end{itemize}
Usa pesos de clase balanceados y una búsqueda en malla (número de árboles, \texttt{max\_features} y \texttt{min\_samples\_leaf}) elegida por macro-F1 de validación.

\textbf{CNN pequeña} (1.2\,M parámetros), entrenada desde cero. Tiene 4 bloques [Conv$3{\times}3$--BN--ReLU]$\times2$ + MaxPool, con anchos 32--256, seguidos de \textit{pooling} global, \textit{dropout} 0.3 y una capa lineal.

\textbf{ResNet-50}~\cite{he2016resnet} (23.5\,M). Cada bloque aprende un residuo, $\mathbf{y}=\mathcal{F}(\mathbf{x},\{W_i\})+\mathbf{x}$; la convolución impone localidad e invariancia a la traslación.

\textbf{ViT-S/16}~\cite{dosovitskiy2021vit} (21.7\,M). Divide la imagen de $224\times224$ en $N=196$ parches de $16\times16$ y los proyecta a $d=384$ dimensiones. Agrega un token \texttt{[CLS]} y embeddings de posición, y aplica 12 bloques de autoatención multi-cabeza (6 cabezas) con
\begin{equation}
\text{Atención}(Q,K,V)=\mathrm{softmax}\!\left(\frac{QK^\top}{\sqrt{d_h}}\right)V,
\end{equation}
donde cada parche puede atender a todos los demás desde la primera capa.

\textbf{Pesos preentrenados.} Ambos modelos grandes provienen de \texttt{timm}~\cite{wightman2019timm}:
\begin{itemize}
\item ResNet-50: \texttt{resnet50.a1\_in1k}, preentrenada en ImageNet-1k.
\item ViT-S/16: \texttt{vit\_small\_patch16\_224.}\allowbreak\texttt{augreg\_in21k\_ft\_in1k}, preentrenado en ImageNet-21k y ajustado en ImageNet-1k.
\end{itemize}
Reemplazamos la cabeza por una capa de 7 salidas y ajustamos \emph{todos} los parámetros. En la ablación ``desde cero'' usamos la misma arquitectura con inicialización aleatoria.

\subsection{Función de pérdida y entrenamiento}
Para compensar el desbalance usamos entropía cruzada ponderada:
\begin{equation}
\mathcal{L}=-\frac{\sum_{i=1}^{B} w_{y_i}\log p_{i,y_i}}{\sum_{i=1}^{B} w_{y_i}},\qquad
p_{i,k}=\frac{e^{z_{i,k}}}{\sum_{j=1}^{K}e^{z_{i,j}}},
\end{equation}
con pesos $w_k=\frac{1/n_k}{\frac{1}{K}\sum_{j}1/n_j}$, donde $n_k$ es el número de imágenes de la clase $k$ en entrenamiento y $K=7$. La pérdida es la media ponderada del minibatch de tamaño $B$, que es lo que calcula \texttt{CrossEntropyLoss} de PyTorch cuando recibe pesos. Así, la clase nematodos (47 imágenes) pesa 3.5 y la clase hongos, 0.32.

\textbf{Optimizador y calendario.} Usamos AdamW~\cite{loshchilov2019adamw} (Tabla~\ref{tab:hp}), con tasas de aprendizaje tomadas de configuraciones estándar de la literatura~\cite{dosovitskiy2021vit,steiner2022howtotrain} y verificadas en validación. La tasa de aprendizaje sube linealmente durante 2 épocas y luego decae con coseno hasta cero:
\begin{equation}
\eta_t=\begin{cases}\eta_{\max}\,\frac{t+1}{T_w}, & t<T_w\\[2pt]
\frac{\eta_{\max}}{2}\Big(1+\cos\frac{\pi(t-T_w)}{T-T_w}\Big), & t\geq T_w\end{cases}
\end{equation}
con $t$ en pasos de minibatch.

\textbf{Configuración común.} Tamaño de lote 32, 40 épocas sin parada temprana, recorte de gradiente a norma 1 y \textit{dropout} en la cabeza: 0.2 en ResNet y ViT, y 0.3 en la CNN pequeña. La precisión mixta fp16 se usa en ResNet y ViT; la CNN pequeña entrena en fp32 porque en fp16 colapsaba en nuestra GPU local.

\textbf{Aumento de datos.} Recorte aleatorio redimensionado (escala 0.6--1), volteos horizontal y vertical, rotación de $\pm20^\circ$ y variación de color: brillo, contraste y saturación ($\pm$20\,\%) y tono ($\pm0.02$). Cada configuración se entrena con las semillas 0, 1 y 2.

\begin{table}[t]
\caption{Hiperparámetros (AdamW, lote 32, 40 épocas, 224$\times$224)}
\label{tab:hp}
\centering\small
\resizebox{\columnwidth}{!}{
\begin{tabular}{lccc}
\toprule
Modelo & $\eta_{\max}$ (pre / cero) & \textit{Weight decay} & Parámetros \\
\midrule
CNN pequeña & -- / $10^{-3}$ & 0.01 & 1.2\,M \\
ResNet-50 & $3{\times}10^{-4}$ / $10^{-3}$ & 0.05 & 23.5\,M \\
ViT-S/16 & $10^{-4}$ / $5{\times}10^{-4}$ & 0.05 & 21.7\,M \\
\bottomrule
\end{tabular}%
}                              
\end{table}

\subsection{Métricas y prueba estadística}
Como las clases están desbalanceadas, la métrica principal es el \textbf{macro-F1}, $\frac{1}{K}\sum_k \mathrm{F1}_k$, que da el mismo peso a cada clase. También reportamos la exactitud balanceada y la exactitud, como media $\pm$ desviación estándar sobre semillas.

Para comparar ResNet-50 y ViT-S/16 con la misma semilla aplicamos la prueba exacta de McNemar~\cite{dietterich1998tests}. Sean $b$ y $c$ las imágenes de prueba que solo acierta uno u otro modelo; bajo $H_0$, $b\sim\mathrm{Binomial}(b+c,\tfrac12)$.

\subsection{Interpretabilidad}
Aplicamos Grad-CAM~\cite{selvaraju2017gradcam}:
\begin{equation}
\alpha_k^c=\frac{1}{Z}\sum_{i,j}\frac{\partial y^c}{\partial A^k_{ij}},\qquad
L^c=\mathrm{ReLU}\Big(\sum_k \alpha_k^c A^k\Big).
\end{equation}
Las capas objetivo son:
\begin{itemize}
\item ResNet-50: el último bloque de \texttt{layer4} (mapa de $7\times7$).
\item CNN pequeña: la última convolución.
\item ViT: la normalización previa al último bloque; descartamos el token \texttt{[CLS]} y reordenamos los 196 tokens en una rejilla de $14\times14$.
\end{itemize}
Para cuantificar la forma de los mapas medimos dos cosas:
\begin{itemize}
\item La \emph{masa top-10\,\%}: fracción de la activación que cae en el 10\,\% de píxeles más activos. Un mapa uniforme da 0.10.
\item La \emph{masa central}: fracción en el cuadrado central, que ocupa el 25\,\% del área. Un mapa uniforme da 0.25.
\end{itemize}

%=====================================================================
\section{Resultados}\label{sec:results}

\begin{table}[t]
\caption{Resultados en prueba (453 imágenes), media $\pm$ DE sobre 3 semillas. \textbf{Negrita}: mejor.}
\label{tab:main}
\centering\small
\setlength{\tabcolsep}{3pt}
\resizebox{\columnwidth}{!}{%
\begin{tabular}{lccc}
\toprule
Modelo & Macro-F1 & Exact. bal. & Exactitud \\
\midrule
Random Forest (color+HOG+LBP) & \pmsd{0.464}{0.021} & \pmsd{0.452}{0.014} & \pmsd{0.558}{0.010} \\
CNN pequeña (cero) & \pmsd{0.614}{0.011} & \pmsd{0.661}{0.004} & \pmsd{0.630}{0.009} \\
ResNet-50 (cero) & \pmsd{0.642}{0.007} & \pmsd{0.687}{0.013} & \pmsd{0.659}{0.003} \\
ViT-S/16 (cero) & \pmsd{0.499}{0.028} & \pmsd{0.550}{0.019} & \pmsd{0.497}{0.029} \\
ResNet-50 (preentrenada) & \pmsd{0.843}{0.012} & \pmsd{0.845}{0.011} & \pmsd{0.855}{0.005} \\
ViT-S/16 (preentrenado) & \textbf{\pmsd{0.853}{0.013}} & \textbf{\pmsd{0.849}{0.008}} & \textbf{\pmsd{0.869}{0.016}} \\
\midrule
\multicolumn{4}{l}{\textit{Publicado~\cite{shabrina2024potato}: partición 90/10 sin deduplicar, 1 corrida}} \\
ResNet-50 & -- & -- & 0.682 \\
EfficientNetV2-B3 & -- & -- & 0.736 \\
\bottomrule
\end{tabular}}
\end{table}

\subsection{Comparación con las líneas base}
La Tabla~\ref{tab:main} resume las 18 corridas propias (6 configuraciones $\times$ 3 semillas).
\begin{itemize}
\item \textbf{Frente a la línea base clásica:} los dos modelos preentrenados superan al Random Forest en más de 0.37 de macro-F1, y a la CNN pequeña en más de 0.22.
\item \textbf{Frente a lo publicado:} nuestra ResNet-50 preentrenada obtiene 0.855 de exactitud, contra 0.682 en~\cite{shabrina2024potato}. Hay que leer esta comparación con cautela, porque la partición y el protocolo difieren.
\item \textbf{Variabilidad entre semillas:} es baja ($\leq0.013$ de DE en macro-F1) para los modelos preentrenados.
\end{itemize}

\subsection{CNN frente a ViT}
Con preentrenamiento, el ViT-S/16 supera a la ResNet-50 por solo 0.010 de macro-F1, una diferencia menor que la desviación entre semillas. La prueba de McNemar por semilla (Tabla~\ref{tab:mcnemar}) no rechaza $H_0$ en ningún caso ($p=0.38$, 1.00 y 0.16).

Por clase (Fig.~\ref{fig:f1}), los dos modelos se reparten las victorias: el ViT es mejor en \textit{Phytophthora} (0.87 frente a 0.80) y hongos (0.85 frente a 0.82); la ResNet, en virus (0.88 frente a 0.86) y nematodos (0.80 frente a 0.77). Concluimos que, preentrenadas, \textbf{no encontramos diferencia significativa entre ambas arquitecturas} en este problema. Esto no prueba que sean equivalentes, ya que para ello haría falta una prueba de equivalencia; pero la diferencia observada (0.010) es menor que la variación entre semillas, así que con estos datos la elección de arquitectura pesa poco.

\begin{table}[t]
\caption{McNemar exacto, ResNet-50 vs ViT-S/16 preentrenados (prueba)}
\label{tab:mcnemar}
\centering\small
\begin{tabular}{cccc}
\toprule
Semilla & Solo ResNet acierta ($b$) & Solo ViT acierta ($c$) & $p$ \\
\midrule
0 & 20 & 27 & 0.38 \\
1 & 24 & 25 & 1.00 \\
2 & 20 & 31 & 0.16 \\
\bottomrule
\end{tabular}
\end{table}

\begin{figure*}[t]
\centering
\includegraphics[width=0.82\textwidth]{figs/f1_por_clase.png}
\caption{F1 por clase en prueba (media y DE sobre 3 semillas; $n$ = imágenes de prueba por clase). Los modelos preentrenados dominan en todas las clases; la mayor brecha con las líneas base está en nematodos, sana y plaga. Las etiquetas se muestran con los nombres originales del dataset en inglés (\textit{fungi}=hongos, \textit{healthy}=sana, \textit{nematode}=nematodos, \textit{pest}=plaga).}
\label{fig:f1}
\end{figure*}

\subsection{Ablación: efecto del preentrenamiento}
La Tabla~\ref{tab:abl} aísla el preentrenamiento manteniendo arquitectura, datos, partición e hiperparámetros, y ajustando solo la tasa de aprendizaje.
\begin{itemize}
\item \textbf{Ganancia en el ViT:} $+0.355$ de macro-F1.
\item \textbf{Ganancia en la ResNet:} $+0.201$.
\item \textbf{ViT desde cero:} queda incluso por debajo de la CNN pequeña, que tiene 18 veces menos parámetros, y es el modelo más inestable entre semillas (DE 0.028).
\item \textbf{ResNet desde cero:} supera levemente a la CNN pequeña (0.642 frente a 0.614).
\end{itemize}

Las curvas de entrenamiento (Fig.~\ref{fig:curvas}) muestran la misma diferencia en la dinámica:
\begin{itemize}
\item \textbf{Preentrenados:} superan 0.8 de macro-F1 de validación en menos de 5 épocas y se estabilizan cerca de 0.9. La pérdida de entrenamiento llega casi a cero, pero la validación no empeora; elegir la mejor época por validación evita usar un punto sobreajustado.
\item \textbf{Desde cero:} a las 40 épocas la pérdida sigue bajando y la validación sigue subiendo. Con más épocas probablemente mejorarían, así que la ablación mide la ganancia del preentrenamiento \emph{con el mismo presupuesto de cómputo}, no su límite.
\end{itemize}

\begin{figure*}[t]
\centering
\includegraphics[width=0.92\textwidth]{figs/curvas_entrenamiento.png}
\caption{Curvas de entrenamiento (media y banda de $\pm1$ DE sobre 3 semillas). (a) Pérdida de entrenamiento. (b) Macro-F1 en validación. Línea continua: preentrenado (y CNN pequeña); discontinua: desde cero. Mismo color = misma arquitectura.}
\label{fig:curvas}
\end{figure*}

Esto coincide con~\cite{dosovitskiy2021vit,touvron2021deit}: sin el sesgo inductivo de localidad, el Transformer necesita aprender desde los datos cómo se relacionan los parches vecinos, y 2\,112 imágenes no alcanzan.

\begin{table}[t]
\caption{Ablación del preentrenamiento (macro-F1 en prueba)}
\label{tab:abl}
\centering\small
\begin{tabular}{lccc}
\toprule
Modelo & Preentrenado & Desde cero & $\Delta$ \\
\midrule
ResNet-50 & \pmsd{0.843}{0.012} & \pmsd{0.642}{0.007} & $+0.201$ \\
ViT-S/16 & \pmsd{0.853}{0.013} & \pmsd{0.499}{0.028} & $+0.355$ \\
\bottomrule
\end{tabular}
\end{table}

\subsection{Análisis de errores}
La Fig.~\ref{fig:conf} muestra la matriz de confusión del ViT preentrenado, sumando las 3 semillas. Los errores se concentran en tres confusiones:
\begin{itemize}
\item \textbf{Plaga $\to$ hongos:} 14\,\% de las fotos de plaga. En ResNet, 11\,\%.
\item \textbf{\textit{Phytophthora} $\to$ hongos:} 12\,\%.
\item \textbf{Nematodos $\to$ hongos:} 17\,\%, con solo 10 imágenes de prueba.
\end{itemize}
Hongos es la clase más numerosa y la de síntomas más variados (manchas necróticas, clorosis), por lo que funciona como ``clase sumidero''. En el Random Forest, el 34\,\% de las fotos de plaga se clasifica como hongos. La propia auditoría apunta a una frontera ambigua: de los 6 pares de fotos idénticas con dos etiquetas, 5 involucran a plaga y 3 son exactamente hongos/plaga (notebook de exploración). Además, el 11\,\% de las hojas sanas se confunde con virus, cuyos síntomas tempranos (mosaico leve) son sutiles a $224\times224$.

\begin{figure}[t]
\centering
\includegraphics[width=0.86\columnwidth]{figs/confusion_vit_small_pre.png}
\caption{Matriz de confusión del ViT-S/16 preentrenado en prueba (\% por fila, 3 semillas sumadas). Hongos actúa como clase sumidero de plaga, \textit{Phytophthora} y nematodos.}
\label{fig:conf}
\end{figure}

\subsection{Interpretabilidad (Grad-CAM)}
La Tabla~\ref{tab:cam} cuantifica los mapas sobre 22 imágenes de prueba (14 aciertos, 2 por clase, y 8 errores elegidos al azar) usando la mejor corrida de cada arquitectura (CNN pequeña s2, ResNet-50 preentrenada s1 y ViT-S/16 preentrenado s2). La Fig.~\ref{fig:cam} ilustra 4 casos representativos de ese subconjunto. Cabe notar que la masa central refleja también que las hojas suelen estar centradas en la foto, y que la ResNet calcula Grad-CAM sobre una rejilla de $7\times7$ frente a $14\times14$ del ViT; al interpolarse a $224\times224$, un mapa más grueso tiende a concentrar mayor masa en el top-10\,\%.
\begin{itemize}
\item \textbf{ResNet-50:} produce mapas compactos (masa top-10\,\% = 0.50) centrados en la lesión. Por ejemplo, señala la mancha necrótica en \textit{Phytophthora} (fila~3) y la perforación en plaga (fila~4). En casos complejos ambos modelos fallan: en la fila~1 la ResNet confunde bacteria con \textit{Phytophthora}, y en la fila~2 tanto ResNet como ViT confunden una lesión avanzada de \textit{Phytophthora} con hongos.
\item \textbf{ViT-S/16:} sus mapas son más dispersos (0.29), casi tan repartidos como un mapa uniforme en la región central (0.27 frente a 0.25). Esto es coherente con su atención global~\cite{raghu2021vitsee}, aunque Grad-CAM es menos fiel en Transformers~\cite{chefer2021transformer}. Aun así, en la fila~3 se concentra en la misma lesión que la ResNet.
\item \textbf{CNN pequeña:} muestra \emph{aprendizaje por atajos}~\cite{geirhos2020shortcut}. En la fila~1 acierta ``bacteria'' mirando el plástico de \textit{mulch} del fondo, no la hoja. En la fila~3 predice ``sana'' porque atiende a las nervaduras en lugar de la lesión.
\end{itemize}
Este hallazgo advierte que, en fotos de campo, un acierto no garantiza que el modelo use la evidencia correcta.

\begin{figure}[t]
\centering
\includegraphics[width=\columnwidth]{figs/gradcam_errores.jpg}
\caption{Grad-CAM en imágenes de prueba (verde = acierto, rojo = error). Fila~1: la CNN pequeña acierta usando el plástico del fondo (atajo). Fila~3: ResNet y ViT localizan la lesión de \textit{Phytophthora}; la CNN pequeña no.}
\label{fig:cam}
\end{figure}

\begin{table}[t]
\caption{Concentración de los mapas Grad-CAM (22 imágenes de prueba)}
\label{tab:cam}
\centering\small
\begin{tabular}{lcc}
\toprule
Modelo & Masa top-10\,\% & Masa central \\
\midrule
CNN pequeña & \pmsd{0.37}{0.18} & \pmsd{0.34}{0.13} \\
ResNet-50 & \pmsd{0.50}{0.11} & \pmsd{0.61}{0.18} \\
ViT-S/16 & \pmsd{0.29}{0.11} & \pmsd{0.27}{0.11} \\
\midrule
Mapa uniforme (referencia) & 0.10 & 0.25 \\
\bottomrule
\end{tabular}
\end{table}

\begin{figure*}[t]
\centering
\includegraphics[width=0.8\textwidth]{figs/app_web.jpg}
\caption{App web con el modelo final (capturas en un celular). Izquierda: parcelas de papa en un mapa satelital; la parcela con un diagnóstico de enfermedad pasa a ``Alerta''. Centro: diagnóstico de una foto de prueba de rancha, con el mapa Grad-CAM y la confianza. Derecha: historial con calendario y el diagnóstico guardado.}
\label{fig:app}
\end{figure*}

\subsection{Costo computacional}
Usamos dos equipos: Google Colab (Tesla T4, 16\,GB) y una laptop con GeForce GTX 1650 Ti de 4\,GB. Las 15 corridas de redes neuronales suman $\approx$5.9\,h de GPU; el Random Forest corre en CPU (Tabla~\ref{tab:cost}).
\begin{itemize}
\item \textbf{En la T4:} una corrida de 40 épocas tarda $\approx$15\,min, con un pico de memoria de 1.5\,GB (ViT) y 1.8\,GB (ResNet).
\item \textbf{En la laptop:} la misma corrida de ResNet tarda 66\,min. En esa GPU, la precisión mixta fp16 produjo colapsos (la validación predecía una sola clase), por lo que la CNN pequeña se entrenó en fp32 y las semillas 1 y 2 de la ResNet preentrenada se entrenaron en la T4.
\item \textbf{En inferencia:} el modelo final clasifica una imagen y calcula su Grad-CAM en $\approx$0.2\,s en CPU.
\end{itemize}

\begin{table}[t]
\caption{Costo de entrenamiento (40 épocas por corrida)}
\label{tab:cost}
\centering\small
\setlength{\tabcolsep}{4pt}
\begin{tabular}{llcc}
\toprule
Modelo & Hardware & min/corrida & Corridas \\
\midrule
Random Forest & CPU & 1.8 & 3 \\
CNN pequeña & GTX 1650 Ti & 40.1 & 3 \\
ResNet-50 (pre) & GTX 1650 Ti / T4 & 65.9 / 15.8 & 1 / 2 \\
ResNet-50 (cero) & T4 & 15.9 & 3 \\
ViT-S/16 (pre) & T4 & 14.4 & 3 \\
ViT-S/16 (cero) & T4 & 15.0 & 3 \\
\midrule
\multicolumn{2}{l}{Total GPU} & \multicolumn{2}{c}{$\approx$5.9\,h} \\
\bottomrule
\end{tabular}
\end{table}

\subsection{Modelo final y demo}
Elegimos como modelo final la corrida con mayor macro-F1 de \emph{validación}, sin mirar la prueba: ViT-S/16 preentrenado con semilla 2. Obtuvo 0.929 en validación y, en prueba, 0.864 de macro-F1 y 0.883 de exactitud.

\textbf{App web.} Publicamos el modelo final en una aplicación web para celular, \url{https://agrotec-papa.vercel.app} (Fig.~\ref{fig:app}). Tiene tres pantallas: parcelas en un mapa satelital, diagnóstico e historial. La interfaz reutiliza el diseño de AGROTEC, un prototipo de asistencia agrícola para Piura desarrollado por uno de los autores; en esta versión, el diagnóstico lo hace el modelo de este trabajo. El modelo corre \emph{en el navegador} con ONNX Runtime Web, sin servidor, y la foto no sale del dispositivo:
\begin{itemize}
\item \textbf{Exportación:} convertimos el ViT a ONNX y cuantizamos sus pesos a int8. El archivo pasa de 89 a 24\,MB y en prueba obtiene 0.864 de macro-F1 y 0.883 de exactitud, igual que el modelo original.
\item \textbf{Grad-CAM sin gradientes:} el navegador no calcula gradientes. Por eso escribimos en el grafo el gradiente analítico del último bloque (cabeza, LayerNorm, MLP y atención del token [CLS]); el mapa coincide con el de PyTorch (error máximo $<10^{-5}$).
\item \textbf{Salida:} la clase más probable con su confianza, las 3 más probables, el mapa Grad-CAM, una recomendación general y una advertencia cuando la confianza es menor que 0.6. Un diagnóstico tarda alrededor de 0.5\,s en una laptop.
\end{itemize}
También hay una demo en Colab y otra local con Gradio, con las mismas salidas.

%=====================================================================
\section{Discusión}\label{sec:discussion}
\textbf{¿Por qué gana el preentrenamiento y no la arquitectura?} Con 2\,112 imágenes, las representaciones aprendidas en ImageNet (bordes, texturas, partes de objetos) son más valiosas que el sesgo inductivo de cada arquitectura. Ese sesgo solo marca diferencia cuando falta el preentrenamiento: la convolución aporta localidad ``gratis'' y por eso la ResNet desde cero aguanta mejor que el ViT.

\textbf{Sesgos y validez.} Hay que tener en cuenta varias limitaciones:
\begin{itemize}
\item \textbf{Generalización geográfica:} todas las fotos provienen de una región de Indonesia, con variedades, iluminación y fondos propios. No sabemos cómo se comportaría el modelo con cultivos andinos.
\item \textbf{Clases pequeñas:} nematodos tiene solo 10 imágenes de prueba, así que su F1 es muy incierto.
\item \textbf{Brecha validación--prueba:} el modelo final obtuvo 0.929 en validación y 0.864 en prueba. Como esta caída de 0.04 a 0.07 aparece en todos los modelos (también en el promedio de corridas no seleccionadas: ViT preentrenado 0.922 en validación frente a 0.853 en prueba; ResNet preentrenada 0.900 frente a 0.843), se debe en parte a que la partición de prueba resultó más difícil que la de validación y en parte a elegir época y modelo con la validación.
\item \textbf{Comparabilidad con lo publicado:} no podemos comparar directamente con~\cite{shabrina2024potato}, porque su partición no controla duplicados.
\item \textbf{Hardware mixto:} la semilla 0 de la ResNet preentrenada se entrenó en la laptop (GTX 1650 Ti) y las semillas 1 y 2 en la Tesla T4 de Colab.
\end{itemize}

\textbf{Uso responsable.} El sistema es una herramienta de apoyo y no reemplaza el diagnóstico de un agrónomo. Una predicción errónea puede llevar a aplicar un fungicida innecesario o a no tratar una plaga a tiempo. Por eso la demo muestra la confianza, el Grad-CAM y advierte cuando la confianza es baja. Antes de usarlo en el Perú haría falta validarlo con fotos locales.

%=====================================================================
\section{Conclusiones}\label{sec:conclusions}
En la detección de 7 enfermedades de la papa con fotos de campo y pocos datos, ResNet-50 y ViT-S/16 preentrenados no muestran diferencia significativa (macro-F1 0.84--0.85) y superan ampliamente a las líneas base. El preentrenamiento es el factor decisivo, en especial para el Transformer ($+0.355$). En los casos analizados, Grad-CAM muestra que la ResNet suele enfocar las lesiones mientras que una red pequeña puede apoyarse en atajos del fondo.

Como trabajo futuro proponemos:
\begin{itemize}
\item Recolectar fotos de papa peruana para evaluar la transferencia de dominio.
\item Probar métodos de explicación específicos para Transformers~\cite{chefer2021transformer}.
\item Destilar el modelo a una versión ligera para celulares.
\end{itemize}

%=====================================================================
\section*{Contribuciones}
\begin{center}\small
\begin{tabular}{>{\raggedright\arraybackslash}p{0.30\columnwidth}>{\raggedright\arraybackslash}p{0.60\columnwidth}}
\toprule
Integrante & Tareas en las que participó \\
\midrule
Gabriel Ángel Nizama Duque & Implementación y entrenamiento de la CNN pequeña, ResNet-50 y ViT-S/16 (3 semillas), ablación del preentrenamiento, selección del modelo final y app web \\
Martín Alexis Valdiviezo Palma & Obtención y auditoría del conjunto de datos, eliminación de duplicados, partición 70/15/15 por grupos, apoyo en la ablación del preentrenamiento y en la app web \\
Juan José Pino Reyes & Línea base Random Forest, prueba de McNemar, Grad-CAM y análisis de errores \\
César Luis Jiménez Miranda & Apoyo en la ejecución de las corridas desde cero en Colab, repositorio y documentación, redacción del artículo y diapositivas \\
\bottomrule
\end{tabular}
\end{center}

\section*{Disponibilidad y uso de herramientas}
El código, la partición, las métricas de las 18 corridas, los pesos del modelo final y los notebooks de reproducción y demo están en \url{https://github.com/gabo2607/papa-cnn-vs-vit}. La app web está en \url{https://agrotec-papa.vercel.app}. Cada tabla y figura se regenera con los comandos del README. Se usó un asistente de IA (Claude) como apoyo en programación y redacción; el equipo revisó y validó el diseño experimental, los resultados y el texto.

\bibliographystyle{IEEEtran}
\bibliography{refs}

\end{document}
